(1) 理想不可压MHD的局域shearing-sheet，背景流 \(\mathbf U_0=-q\Omega x\,\hat y\)，\(0<q<2\)、\(\Omega>0\)，竖直均匀背景磁场 \(B_0\hat z\)。忽略背景密度梯度，取轴对称扰动 \(\exp(st+ikz)\)，且只有 \(x,y\) 分量。把磁扰动按 \(\boldsymbol b=\delta\mathbf B/\sqrt{\mu_0\rho}\) 规范化，推导四条线性化速度—磁场方程。
(2) 消去四个振幅求 \(s\) 的四次色散方程，并以 \(\kappa^2=2(2-q)\Omega^2\)、\(\omega_A=kv_A\) 表示，\(v_A=B_0/\sqrt{\mu_0\rho}\)。核验零磁场时的回转稳定极限。
(3) 给出轴对称MRI增长的准确磁场—波数范围，求一般 \(0<q<2\) 的最快增长率和对应 \(k_{\max}\)。代入 Kepler 剪切 \(q=3/2\)，并说明强磁场为何抑制此局域模式。
(4) 对一组复振幅定义正定扰动能量 \(\mathcal E=\rho(|u_x|^2+|u_y|^2+|b_x|^2+|b_y|^2)/2\)，计算其随时间的变化，明确背景剪切供能对应的 Reynolds 与 Maxwell 应力。检验 \(q=0\) 时的能量守恒，并陈述局域、不可压和理想导体假设的适用边界。
